\section{Глава~\ref{sec:glutcomputer}. Что вычисляют \nt{GLUT}-нейроны}

Логические утверждения главы --- теоремы о схемах из неотрицательных весов, выведенные в самой главе; опыт подтверждает модель нейрона и то, что схемы, собранные по этим теоремам (детектор совпадений, линия задержки, самоподдерживающаяся группа, цепочка), в мозге действительно встречаются.

{\footnotesize
\setlength{\tabcolsep}{3pt}\renewcommand{\arraystretch}{1.15}
\begin{longtable}{>{\raggedright}p{0.5cm}>{\raggedright}p{4.0cm}>{\raggedright}p{5.6cm}>{\raggedright}p{2.3cm}>{\raggedright\arraybackslash}p{1.7cm}}
\toprule
№ & Утверждение & Опыт и что измерено & Источник & Статус \\
\midrule\endfirsthead
\toprule
№ & Утверждение & Опыт и что измерено & Источник & Статус \\
\midrule\endhead
1 & Нейрон --- RC-цепь с порогом и сбросом~\eqref{eq:glutneuron} & В тело пирамидных нейронов коры крысы (срезы) подавали шумовой ток, похожий на синаптический поток в живом мозге. Зависимость средней частоты от среднего и разброса тока описывается интегрирующим нейроном с утечкой, если добавить медленную адаптацию частоты. Диапазоны $\tau$ и задержек приведены в главе без ссылки на опыт & \cite{Rauch2003} & измерено \\
2 & Модель~\eqref{eq:glutneuron} --- упрощение: ветви входов сами выдают местные импульсы & Запись прямо с дендритов пирамидных нейронов зрительной коры мыши in vivo: зрительный раздражитель вызывает местные дендритные импульсы, зависящие от рецепторов NMDA; их подавление расширяет настройку нейрона на ориентацию & \cite{Smith2013} & измерено \\
3 & Окно совпадения $\Delta_{\max}=\tau\ln\frac{w}{\theta-w}$~\eqref{eq:window}; при $\theta=1.5w$ оно равно $0.7\tau$ & Следствие модели~\eqref{eq:glutneuron}. Прямое измерение узкого окна суммирования есть для нейронов с быстрым торможением (глава~\ref{sec:gaba}, строка~3 её таблицы) & вывод в главе & теория \\
4 & Сеть из одних \nt{GLUT}-нейронов вычисляет только монотонные функции входов (утверждение~\ref{prop:monotone}) & Доказательство в главе: итерация от тишины с неубывающей правой частью. Это утверждение о модели; опыта, который проверял бы его на живой сети без торможения, нет & вывод в главе & теория \\
5 & Органы чувств дают пару проводов: одни клетки отвечают на усиление раздражителя, другие --- на ослабление & Сетчатка: уже на биполярных клетках сигнал колбочек расходится на включающий и выключающий каналы. Фармакологический блок включающих биполярных клеток у обезьяны: каналы идут раздельно до коры; после блока ухудшается обнаружение усиления света, но не ослабления & \cite{Schiller1992} & измерено \\
6 & С двухпроводным кодом сеть из \nt{GLUT}-нейронов вычисляет любую логическую функцию асинхронно, без общего такта, ценой вдвое большего числа нейронов (утверждение~\ref{prop:dualrail}, \eqref{eq:dualrail}, рис.~\ref{fig:dualxor}) & Двухпроводный код и определение готовности по самому сигналу --- стандартная техника асинхронных схем, нечувствительных к задержкам. Схема исключающего ИЛИ из шести нейронов собрана в главе. Опыта, показывающего, что мозг вычисляет логические функции именно двухпроводным кодом, нет & \cite{Sparso2001} & теория \\
7 & Наибольшая разница времени прихода звука в уши у человека $\approx0.6\ms$, а мозг различает около десяти микросекунд & Слушатели в опыте с двумя вариантами ответа: порог различения разницы времени (75\,\% правильных) --- $9$\,мкс для шума в полосе $150$--$1700$\,Гц, $11$\,мкс для тона $1$\,кГц, $28$\,мкс для щелчка. Предел $0.6\ms$ --- расчёт главы, $0.22\,\text{м}/343\,\text{м/с}$ & \cite{KlumppEady1956} & измерено \\
8 & У птиц разница времён превращается в место через линии задержки и детекторы совпадений~\eqref{eq:itd} (рис.~\ref{fig:itd}) & Сипуха: аксоны от двух ушей входят в ядро детекторов совпадений с противоположных сторон; время прихода импульсов вдоль аксона меняется упорядоченно с глубиной, и разброс задержек через ядро ($160$\,мкс на $700$\,мкм) совпадает с диапазоном межушных задержек & \cite{Carr1990} & измерено \\
9 & У млекопитающих ряда задержек не нашли; ту же задачу решают детекторы совпадений с точно рассчитанным торможением от \nt{GLYC}-нейронов & Запись in vivo в медиальной верхней оливе песчанки: ответы не образуют карты направлений; подведение глицина и его блокатора через микропипетку показывает, что точно рассчитанное во времени \emph{глициновое} торможение задаёт рабочий диапазон задержек. Торможение здесь от \nt{GLYC}, а не от \nt{GABA} & \cite{Brand2002,Grothe2010} & измерено \\
10 & Группа с сильной обратной связью --- триггер; самоподдерживающаяся активность держится секунды после раздражителя (рис.~\ref{fig:bistable}) & Префронтальная кора обезьяны в задаче с отложенным движением глаз к запомненной точке (задержка $1$--$6$\,с): $87$ из $170$ нейронов, связанных с задачей, меняют частоту во время задержки, у $79\,\%$ из них это зависит от направления на запомненную точку. Что эту активность держит обратная связь с двумя устойчивыми состояниями, --- теория & \cite{Funahashi1989,Wang2001} & измерено \\
11 & Бит записывается, если приток длится дольше $\approx0.7\tau$ (рис.~\ref{fig:recurrentgroup}б) & Расчёт уравнения $\tau\,\dd r/\dd t=-r+f(wr+I)$ методом Эйлера при $w=1$, $I=0.6$; скрипта в \texttt{models/} нет. Опыта нет & расчёт в главе & модель \\
12 & Память можно хранить в облегчении синапсов, почти без импульсов & Теория: остаточный кальций в окончаниях держит запись около секунды. Опора: в медиальной префронтальной коре хорька (запись с нескольких нейронов сразу) есть подсеть пирамидных нейронов, связанных облегчающимися синапсами с долгим последействием. Обзор наблюдений «тихих» промежутков при удержании в памяти. Какой способ кора использует когда --- открытый вопрос & \cite{Mongillo2008,WangMarkram2006,Stokes2015} & гипотеза \\
13 & Память стирается усталостью: запас медиатора истощается & Парные записи пирамидных нейронов коры: ответ синапса падает при частых импульсах, скорость падения задаётся вероятностью выброса. Что именно это гасит самоподдерживающуюся группу, прямо не показано & \cite{Tsodyks1997} & измерено \\
14 & Цепочка групп --- таймер, запускаемый событием; у певчих птиц нейроны срабатывают строго по очереди & Запись опознанных нейронов высшего вокального центра зебровой амадины во сне и при пении: каждый нейрон, посылающий выход в моторное ядро, выдаёт одну короткую пачку в один точный момент песни, и нейроны срабатывают последовательно & \cite{Hahnloser2002} & измерено \\
\bottomrule
\end{longtable}}
